
\documentclass[../sablona-main.tex]{subfiles}
\begin{document}

\section{Identifikace náhodných parametrů (úkol 10)}
V našem modelu předpokládáme, že všechny parametry jsou 100% jisté a fixní. V reálném výrobním procesu to tak ale nikdy nebude a nějaké parametry budou méně stabilní a náhnodné.

Například čas potřebný na montáž (0.6 h a 1.5 h), nikdy nebude přesně 1,5 hodiny, lidský faktor to ovlivní natolik aby nějaký stůl byl rychleji složen a nějaký pomaleji. Nehledě na to že si pracovník může zajít v během skládání na kafe, nebo prohodit pár slov s kolegou.
Také, časové kapacity se mohou lišit, onemocnění pracovníka, porucha vybavení a nebo jiné události mohou zapříčinit snížení kapacit.
Dalším náhodným faktorem je kvalita materiálu. Né každá skleněná deska bude perfektní a v případě poškrábání nebo případného zničení při práci se spálí jak materiál tak časový prostor.
V neposlední řadě model předpokládá že se vše prodá, a to za stejnou cenu, na reálném trhu je však ale situace jiná, a kolísá jak cena tak nabídka a poptávka.
Z tohoto důvodu pracujeme s ideálním modelem pro naši optimalizaci.

\end{document}